% Part: applied-modal-logics
% Chapter: epistemic-logic
% Section: properties-accessibility

\documentclass[../../../include/open-logic-section]{subfiles}

\begin{document}

\olfileid{aml}{el}{acc}

\olsection{د لاسرسي اړيکې او معرفتي اصول}

د نورمال مودالي منطقونو د چوکاټ د مطابقت په اړه له خپلو پخوانيو
معلوماتو سره، ښايي وغواړو وګورو چې ځانګړي !!{formula}s د معرفتي
تفسير له مخې څۀ بڼه لري۔ مخکې مو ويلي چې معرفتي منطقونه عموماً
په S5 کښې تفسيرېږي۔ نو وګورو چې بېلابېل معرفتي اصول څنګه ښودل
کېږي او له بېلابېلو چوکاټي شرطونو سره يې مطابقت څنګه دے۔

له نورمال مودالي منطق څخه ياد کړئ چې بېلابېلو مودالي !!{formula}s
د لاسرسي د اړيکو بېلابېل خاصيتونه مشخصول۔ دا جدول له هغو څخه څو
خاصيتونه رااخلي چې له ټاکلو معرفتي اصولو سره تړاو لري۔

\begin{table}[t]
    \begin{tabular}{| p{.48\textwidth} || p{.48\textwidth} |}
      \hline
      {\emph{که $R$ \dots وي}} & {\emph{نو \dots په~$\mModel{M}$ کښې رښتيا دے:}} \\
      \hline \hline

      & $\Knows (p \lif q) \lif (\Knows p \lif \Knows q)$
      \hfill \newline{(تړلتيا)} \\
      \hline
      \emph{انعکاسي}: $\forall w Rww$
      & $\Knows p \lif p$ \hfill \newline{(رښتينولي)} \\
      \hline
      \emph{متعدي}: \newline
      $\forall u \forall v \forall w ((Ruv \land Rvw) \lif Ruw)$ &
      $\Knows p \lif \Knows \Knows p$ \hfill
           \newline{(مثبت ځان‌پوهاوی)} \\
      \hline
      \emph{اقليدسي}: \newline
      $\forall w \forall u \forall v ((Rwu \land Rwv) \lif Ruv)$ &
      $\lnot \Knows p \lif \Knows \neg \Knows p$ \hfill
        \newline{(منفي ځان‌پوهاوی)} \\
      \hline
    \end{tabular}
    \caption{څلور معرفتي اصول۔}
    \ollabel{tab:four}
  \end{table}

رښتينولي، چې د $T$ اکسيوم سره مطابقت لري، ډېر ځله د دغو اصولو
تر ټولو لږ د اختلاف وړ ګڼل کېږي؛ ځکه دا څرګندوي چې که يو
!!a{formula} پېژندل شوے وي، نو بايد رښتيا هم وي۔ تړلتيا او مثبت
او منفي ځان‌پوهاوی تر دې ډېر د بحث وړ دي۔

تړلتيا، چې د $K$ اکسيوم سره مطابقت لري، دا نظر ښيي چې د يوه
عامل پوهه د استلزام له مخې تړلې ده۔ په ځينو حالاتو کښې ښايي دا
موږ ته معقوله ښکاره شي۔ د بېلګې په توګه، زه ښايي پوهېږم چې که
په ويکټوريا کښې يم، نو د وېنکوور پر ټاپو يم۔ له عجيبو شکاکانه
حالتونو پرته، زه پوهېږم چې په ويکټوريا کښې يم؛ نو بايد دا هم
راوښيي چې زه د وېنکوور پر ټاپو يم۔ په دې حالت کښې زما د پوهې
منطقي تړلتيا ښايي تر يوه بريده طبيعي ښکاره شي۔ خو موږ تل د
خپلې پوهې په ټولو پايلو فکر نۀ کوو، او له دې امله په نورو
حالتونو کښې لږ طبيعي نتيجې هم راوتلے شي۔

مثبت ځان‌پوهاوی، چې کله د KK-اصل هم بلل کېږي، داسې بيانېږي: که
زه پر يوه څۀ پوهېږم، نو پوهېږم چې پرې پوهېږم۔ دا د 4 اکسيوم
معرفتي سيال دے۔ په ورته ډول، منفي ځان‌پوهاوی داسې بيانېږي: که
زه پر يوه څۀ \emph{نۀ} پوهېږم، نو پوهېږم چې پرې نۀ پوهېږم؛ دا د
5 اکسيوم سيال دے۔ دواړو ته نسبتاً عادي ضدبېلګې پيدا کېدے شي:
زه ښايي په دې نۀ يم ډاډه چې آيا پر هغه څۀ پوهېږم که نۀ، چې په
حقيقت کښې پرې پوهېږم۔


\end{document}
